% Timing: 10:00--13:00   (paper: Sec. 1.2)
\section{Natural approaches that fail}

% ============================================================
\begin{frame}{$\Sigma_1$\CoR{}: a fooling-set argument \ \NOTE{\fcite{nakamura2020}, \fcite{nakamura2022}}}

\begin{center}
\begin{tikzpicture}[every node/.style={font=\small}]
  % small structures
  \node[structbox,minimum width=2.6cm,minimum height=1.3cm] (Si) at (0,1.4) {};
  \node[above] at (Si.north) {$\mathcal S_i$};
  \node[pv,label=left:$u_i$] (ui) at (-.9,1.2) {};
  \node[pv,label=above:$w_i$] (wi) at (0,1.75) {};
  \node[pv,label=right:$v_i$] (vi) at (.9,1.2) {};
  \draw[thick,-{Stealth}] (ui) -- (wi); \draw[thick,-{Stealth}] (wi) -- (vi);
  \node[structbox,minimum width=2.6cm,minimum height=1.3cm] (Sj) at (0,-1.4) {};
  \node[below] at (Sj.south) {$\mathcal S_j$};
  \node[pv,label=left:$u_j$] (uj) at (-.9,-1.6) {};
  \node[pv,label=below:$w_j$] (wj) at (0,-1.05) {};
  \node[pv,label=right:$v_j$] (vj) at (.9,-1.6) {};
  \draw[thick,-{Stealth}] (uj) -- (wj); \draw[thick,-{Stealth}] (wj) -- (vj);
  \node[font=\scriptsize,text=chibasub] at (2.0,1.4) {$\varphi = 1$};
  \node[font=\scriptsize,text=chibasub] at (2.0,-1.4) {$\varphi = 1$};
  % large structure
  \begin{scope}[visible on=<2->]
  \node[structbox,minimum width=4.4cm,minimum height=2.4cm] (N) at (6.4,0) {};
  \node[above] at (N.north) {$\mathcal N$};
  \node[pv,label=left:$u^*$] (u) at (4.8,-.3) {};
  \node[pv,label=above:$c$] (c) at (6.4,.6) {};
  \node[pv,label=right:$v^*$] (v) at (8.0,-.3) {};
  \draw[thick,-{Stealth}] (u) -- (c); \draw[thick,-{Stealth}] (c) -- (v);
  \node[font=\scriptsize,text=chibasub] at (6.4,-.85) {$\varphi = 0$};
  \draw[-{Stealth},densely dashed,black!50] (Si.east) to[bend left=15] node[above,font=\scriptsize] {$\hookrightarrow$} (5.5,.4);
  \draw[-{Stealth},densely dashed,black!50] (Sj.east) to[bend right=15] node[below,font=\scriptsize] {$\hookrightarrow$} (7.3,.1);
  \end{scope}
\end{tikzpicture}
\end{center}

\vfill
\onslide<3->{%
\begin{center}
  $\Sigma_1$\CoR: \ $\iota(\sem{T}{\mathcal S}) \subseteq \sem{T}{\mathcal S'}$ \ $(\diamond)$
  \quad $\Rightarrow$ \quad each composition separates $\le 1$ instance
\end{center}}

\note{
\textbf{伝えること}：既存の指数下界は $\Sigma_1$\CoR{}（補が合成の上に来ない）だけ。その証明は fooling-set 論法として整理できる。
\begin{itemize}
  \item 小さい構造 $\mathcal S_i$ は $\varphi = 1$、大きい構造 $\mathcal N$ は $\varphi = 0$、交差した組は $\mathcal N$ に埋め込める
  \item 拡張で保存される $(\diamond)$ ので、各合成は 1 つのインスタンスしか分離できない $\to$ $t$ 個必要
\end{itemize}
}
\end{frame}

% ============================================================
\begin{frame}{With complement, both ingredients fail}

\begin{center}
\begin{tikzpicture}[every node/.style={font=\small}]
  \node[structbox,minimum width=2.6cm,minimum height=1.6cm] (S) at (0,0) {};
  \node[above] at (S.north) {$\mathcal S$};
  \node[pv,label=left:$u$] (u) at (-.5,.1) {};
  \node[font=\scriptsize] at (0,-.5) {$\overline{R \mathbin{;} \one}$ holds};
  \draw[-{Stealth},thick] (1.7,0) -- node[above] {$\iota$} (2.9,0);
  \node[structbox,minimum width=2.6cm,minimum height=1.6cm] (S2) at (4.6,0) {};
  \node[above] at (S2.north) {$\mathcal S'$};
  \node[pv,label=left:$\iota(u)$] (u2) at (4.2,.1) {};
  \node[pv] (r) at (5.3,.45) {};
  \draw[thick,-{Stealth},chibasub] (u2) -- node[above,font=\scriptsize] {$R$} (r);
  \node[font=\scriptsize,text=chibasub] at (4.6,-.5) {$\overline{R \mathbin{;} \one}$ fails};
\end{tikzpicture}

1. complement reverses $(\diamond)$
\end{center}

\vfill
\pause
\begin{center}
  2. \ $\overline{(E \cap (E \mathbin{;} \overline{\id}))} \mathbin{;} \one$ \ separates every $\mathcal S_i$ from $\mathcal N$
  \ \NOTE{(two compositions)}
\end{center}

\vfill
\pause
\HP{Unlike the inclusion $(\diamond)$, an \EMPH{equality} survives complement.}

\note{
\textbf{伝えること}：補を許すと、この論法の 2 つの材料が両方壊れる。教訓：\emph{包含}は補で反転するが、\emph{等式}は補を通る。
\begin{itemize}
  \item 1. 補は $(\diamond)$ を逆向きにする（図）
  \item 2. 合成 2 つの項で全インスタンスを分離できてしまう
  \item 他の手法（formula size、EF ゲーム、rank）も直接は使えない
  \item 最後の一行が、我々の不変条件を等式にする動機
\end{itemize}
}
\end{frame}
